\section{Conclusion}
\label{sec:conclusion}

In this paper, we propose \EdgeQuantum, a scalable tiered-memory quantum simulation framework for resource-constrained edge devices. By coordinating UMA-aware zero-copy data movement with asynchronous I/O serialization and mantissa-truncated LZ4 ping-pong compression, we mask NVMe bottlenecks and extend simulation capacity beyond physical DRAM. We evaluate \EdgeQuantum on an 8~GB NVIDIA Jetson Orin Nano, achieving up to 4.59$\times$ speedup over blocking I/O baselines and supporting simulation up to 39 qubits (4~TB), a 512$\times$ logical capacity expansion.

\begin{comment}
 In this paper, we observe that existing quantum circuit simulators on edge devices are severely limited by physical memory capacity and the high latency of commodity storage, making large-scale simulation impossible.
To address these challenges, we propose \EdgeQuantum, a scalable tiered-memory simulation framework that treats NVMe storage as a seamless extension of GPU memory.
\EdgeQuantum integrates a novel 4-buffer asynchronous pipeline to serialize read/write operations and hide NVMe contention, and employs an adaptive mantissa-truncated LZ4 ping-pong compression strategy to maximize effective storage capacity.
Our evaluations show that \EdgeQuantum enables the simulation of up to 39 qubits (4~TB state vector) on an 8~GB NVIDIA Jetson Orin Nano, achieving a 512$\times$ memory capacity expansion beyond the physical RAM limit.
Furthermore, it outperforms blocking I/O baselines by up to 1.48$\times$ and successfully executes complex circuits such as 39-qubit Random and QV, which were previously exclusive to HPC-class infrastructure.
These results demonstrate that \EdgeQuantum effectively breaks the memory wall of edge computing, offering a practical solution for ubiquitous and decentralized quantum algorithm development.   
\end{comment}
